\documentclass[12pt,a4paper]{article}
\usepackage[top=1.8cm,bottom=1.8cm,left=1.8cm,right=1.8cm]{geometry}
\usepackage{fontspec}
\setmainfont{TeX Gyre Termes}
\usepackage{amsmath,amssymb}
\usepackage{fancyhdr}
\usepackage{tasks}
\usepackage{tikz}
\settasks{label=\textbf{\Alph*.},label-width=1.5em,item-indent=2em}
\setlength{\parindent}{0pt}
\setlength{\parskip}{3pt}
\pagestyle{fancy}\fancyhf{}\renewcommand{\headrulewidth}{0pt}
\fancyhead[L]{\small Ngân hàng bài tập Toán 9 -- 9H1-C1-B1-D02}
\cfoot{\small Trang \thepage}
\newif\ifloigiai
\loigiaitrue
\newenvironment{cauhoi}[2]{%
  \par\medskip\noindent\rule{\textwidth}{0.4pt}\par\nopagebreak
  \noindent\textbf{[#1]}\hfill{\scriptsize\ttfamily #2}\par\nopagebreak\smallskip}{\par}
\newcommand{\loigiai}[1]{\ifloigiai\par\smallskip\noindent\textit{Lời giải.} #1\par\fi}
\newcommand{\dapan}[1]{\ifloigiai\par\smallskip\noindent\textit{Đáp án:} \textbf{#1}\par\fi}
\begin{document}
\begin{center}
{\Large\textbf{NGÂN HÀNG BÀI TẬP TOÁN 9}}\\[3pt]
\textbf{Chương 1. Phương trình và hệ hai phương trình bậc nhất hai ẩn}\\[2pt]
\textbf{Bài 1. Khái niệm phương trình và hệ hai phương trình bậc nhất hai ẩn}\\[2pt]
\textbf{Dạng D02: Nghiệm tổng quát}\\[2pt]
\textit{Mã dạng: 9H1-C1-B1-D02 \quad 50 bài tự luận + 50 câu trắc nghiệm \quad (TH: 50)}
\end{center}
\textbf{Quy ước mã câu:} \textit{khối, học kì, chương, bài, dạng, loại (TL/TN), mức độ, số thứ tự}.\par Ví dụ: \texttt{9H1-C1-B1-D02-TL-NB-001}.\par Để in bản cho học sinh, đặt \verb|\loigiaifalse| ngay sau \verb|\loigiaitrue|.
\section*{Phần I. Tự luận}
\begin{cauhoi}{Bài 1}{9H1-C1-B1-D02-TL-TH-001}
Viết nghiệm tổng quát của phương trình $-x+5y=-9$ (biểu diễn $y$ theo $x$).
\loigiai{Ta có $5y=-9+x$ nên $y=\dfrac{1}{5}x-\dfrac{9}{5}$. Nghiệm tổng quát là $(x;\,\dfrac{1}{5}x-\dfrac{9}{5})$ với $x\in\mathbb R$.}
\end{cauhoi}
\begin{cauhoi}{Bài 2}{9H1-C1-B1-D02-TL-TH-002}
Viết nghiệm tổng quát của phương trình $-4x+3y=-5$ (biểu diễn $y$ theo $x$).
\loigiai{Ta có $3y=-5+4x$ nên $y=\dfrac{4}{3}x-\dfrac{5}{3}$. Nghiệm tổng quát là $(x;\,\dfrac{4}{3}x-\dfrac{5}{3})$ với $x\in\mathbb R$.}
\end{cauhoi}
\begin{cauhoi}{Bài 3}{9H1-C1-B1-D02-TL-TH-003}
Viết nghiệm tổng quát của phương trình $2x-3y=6$ (biểu diễn $y$ theo $x$).
\loigiai{Ta có $-3y=6-2x$ nên $y=\dfrac{2}{3}x-2$. Nghiệm tổng quát là $(x;\,\dfrac{2}{3}x-2)$ với $x\in\mathbb R$.}
\end{cauhoi}
\begin{cauhoi}{Bài 4}{9H1-C1-B1-D02-TL-TH-004}
Viết nghiệm tổng quát của phương trình $-5x+2y=-1$ (biểu diễn $y$ theo $x$).
\loigiai{Ta có $2y=-1+5x$ nên $y=\dfrac{5}{2}x-\dfrac{1}{2}$. Nghiệm tổng quát là $(x;\,\dfrac{5}{2}x-\dfrac{1}{2})$ với $x\in\mathbb R$.}
\end{cauhoi}
\begin{cauhoi}{Bài 5}{9H1-C1-B1-D02-TL-TH-005}
Viết nghiệm tổng quát của phương trình $3x+5y=3$ (biểu diễn $y$ theo $x$).
\loigiai{Ta có $5y=3-3x$ nên $y=-\dfrac{3}{5}x+\dfrac{3}{5}$. Nghiệm tổng quát là $(x;\,-\dfrac{3}{5}x+\dfrac{3}{5})$ với $x\in\mathbb R$.}
\end{cauhoi}
\begin{cauhoi}{Bài 6}{9H1-C1-B1-D02-TL-TH-006}
Viết nghiệm tổng quát của phương trình $4x+3y=7$ (biểu diễn $y$ theo $x$).
\loigiai{Ta có $3y=7-4x$ nên $y=-\dfrac{4}{3}x+\dfrac{7}{3}$. Nghiệm tổng quát là $(x;\,-\dfrac{4}{3}x+\dfrac{7}{3})$ với $x\in\mathbb R$.}
\end{cauhoi}
\begin{cauhoi}{Bài 7}{9H1-C1-B1-D02-TL-TH-007}
Viết nghiệm tổng quát của phương trình $-2x-3y=-4$ (biểu diễn $y$ theo $x$).
\loigiai{Ta có $-3y=-4+2x$ nên $y=-\dfrac{2}{3}x+\dfrac{4}{3}$. Nghiệm tổng quát là $(x;\,-\dfrac{2}{3}x+\dfrac{4}{3})$ với $x\in\mathbb R$.}
\end{cauhoi}
\begin{cauhoi}{Bài 8}{9H1-C1-B1-D02-TL-TH-008}
Viết nghiệm tổng quát của phương trình $-3x+5y=4$ (biểu diễn $y$ theo $x$).
\loigiai{Ta có $5y=4+3x$ nên $y=\dfrac{3}{5}x+\dfrac{4}{5}$. Nghiệm tổng quát là $(x;\,\dfrac{3}{5}x+\dfrac{4}{5})$ với $x\in\mathbb R$.}
\end{cauhoi}
\begin{cauhoi}{Bài 9}{9H1-C1-B1-D02-TL-TH-009}
Viết nghiệm tổng quát của phương trình $2x+5y=-5$ (biểu diễn $y$ theo $x$).
\loigiai{Ta có $5y=-5-2x$ nên $y=-\dfrac{2}{5}x-1$. Nghiệm tổng quát là $(x;\,-\dfrac{2}{5}x-1)$ với $x\in\mathbb R$.}
\end{cauhoi}
\begin{cauhoi}{Bài 10}{9H1-C1-B1-D02-TL-TH-010}
Viết nghiệm tổng quát của phương trình $2x+5y=-6$ (biểu diễn $y$ theo $x$).
\loigiai{Ta có $5y=-6-2x$ nên $y=-\dfrac{2}{5}x-\dfrac{6}{5}$. Nghiệm tổng quát là $(x;\,-\dfrac{2}{5}x-\dfrac{6}{5})$ với $x\in\mathbb R$.}
\end{cauhoi}
\begin{cauhoi}{Bài 11}{9H1-C1-B1-D02-TL-TH-011}
Viết nghiệm tổng quát của phương trình $-3x+4y=5$ (biểu diễn $y$ theo $x$).
\loigiai{Ta có $4y=5+3x$ nên $y=\dfrac{3}{4}x+\dfrac{5}{4}$. Nghiệm tổng quát là $(x;\,\dfrac{3}{4}x+\dfrac{5}{4})$ với $x\in\mathbb R$.}
\end{cauhoi}
\begin{cauhoi}{Bài 12}{9H1-C1-B1-D02-TL-TH-012}
Viết nghiệm tổng quát của phương trình $4x+5y=5$ (biểu diễn $y$ theo $x$).
\loigiai{Ta có $5y=5-4x$ nên $y=-\dfrac{4}{5}x+1$. Nghiệm tổng quát là $(x;\,-\dfrac{4}{5}x+1)$ với $x\in\mathbb R$.}
\end{cauhoi}
\begin{cauhoi}{Bài 13}{9H1-C1-B1-D02-TL-TH-013}
Viết nghiệm tổng quát của phương trình $2x-3y=-1$ (biểu diễn $y$ theo $x$).
\loigiai{Ta có $-3y=-1-2x$ nên $y=\dfrac{2}{3}x+\dfrac{1}{3}$. Nghiệm tổng quát là $(x;\,\dfrac{2}{3}x+\dfrac{1}{3})$ với $x\in\mathbb R$.}
\end{cauhoi}
\begin{cauhoi}{Bài 14}{9H1-C1-B1-D02-TL-TH-014}
Viết nghiệm tổng quát của phương trình $x+2y=4$ (biểu diễn $y$ theo $x$).
\loigiai{Ta có $2y=4-x$ nên $y=-\dfrac{1}{2}x+2$. Nghiệm tổng quát là $(x;\,-\dfrac{1}{2}x+2)$ với $x\in\mathbb R$.}
\end{cauhoi}
\begin{cauhoi}{Bài 15}{9H1-C1-B1-D02-TL-TH-015}
Viết nghiệm tổng quát của phương trình $-x+4y=-3$ (biểu diễn $y$ theo $x$).
\loigiai{Ta có $4y=-3+x$ nên $y=\dfrac{1}{4}x-\dfrac{3}{4}$. Nghiệm tổng quát là $(x;\,\dfrac{1}{4}x-\dfrac{3}{4})$ với $x\in\mathbb R$.}
\end{cauhoi}
\begin{cauhoi}{Bài 16}{9H1-C1-B1-D02-TL-TH-016}
Viết nghiệm tổng quát của phương trình $2x-3y=8$ (biểu diễn $y$ theo $x$).
\loigiai{Ta có $-3y=8-2x$ nên $y=\dfrac{2}{3}x-\dfrac{8}{3}$. Nghiệm tổng quát là $(x;\,\dfrac{2}{3}x-\dfrac{8}{3})$ với $x\in\mathbb R$.}
\end{cauhoi}
\begin{cauhoi}{Bài 17}{9H1-C1-B1-D02-TL-TH-017}
Viết nghiệm tổng quát của phương trình $-x+5y=-1$ (biểu diễn $y$ theo $x$).
\loigiai{Ta có $5y=-1+x$ nên $y=\dfrac{1}{5}x-\dfrac{1}{5}$. Nghiệm tổng quát là $(x;\,\dfrac{1}{5}x-\dfrac{1}{5})$ với $x\in\mathbb R$.}
\end{cauhoi}
\begin{cauhoi}{Bài 18}{9H1-C1-B1-D02-TL-TH-018}
Viết nghiệm tổng quát của phương trình $5x+4y=8$ (biểu diễn $y$ theo $x$).
\loigiai{Ta có $4y=8-5x$ nên $y=-\dfrac{5}{4}x+2$. Nghiệm tổng quát là $(x;\,-\dfrac{5}{4}x+2)$ với $x\in\mathbb R$.}
\end{cauhoi}
\begin{cauhoi}{Bài 19}{9H1-C1-B1-D02-TL-TH-019}
Viết nghiệm tổng quát của phương trình $-x+2y=5$ (biểu diễn $y$ theo $x$).
\loigiai{Ta có $2y=5+x$ nên $y=\dfrac{1}{2}x+\dfrac{5}{2}$. Nghiệm tổng quát là $(x;\,\dfrac{1}{2}x+\dfrac{5}{2})$ với $x\in\mathbb R$.}
\end{cauhoi}
\begin{cauhoi}{Bài 20}{9H1-C1-B1-D02-TL-TH-020}
Viết nghiệm tổng quát của phương trình $-2x+4y=-7$ (biểu diễn $y$ theo $x$).
\loigiai{Ta có $4y=-7+2x$ nên $y=\dfrac{1}{2}x-\dfrac{7}{4}$. Nghiệm tổng quát là $(x;\,\dfrac{1}{2}x-\dfrac{7}{4})$ với $x\in\mathbb R$.}
\end{cauhoi}
\begin{cauhoi}{Bài 21}{9H1-C1-B1-D02-TL-TH-021}
Viết nghiệm tổng quát của phương trình $4x-2y=1$ (biểu diễn $y$ theo $x$).
\loigiai{Ta có $-2y=1-4x$ nên $y=2x-\dfrac{1}{2}$. Nghiệm tổng quát là $(x;\,2x-\dfrac{1}{2})$ với $x\in\mathbb R$.}
\end{cauhoi}
\begin{cauhoi}{Bài 22}{9H1-C1-B1-D02-TL-TH-022}
Viết nghiệm tổng quát của phương trình $3x+2y=-1$ (biểu diễn $y$ theo $x$).
\loigiai{Ta có $2y=-1-3x$ nên $y=-\dfrac{3}{2}x-\dfrac{1}{2}$. Nghiệm tổng quát là $(x;\,-\dfrac{3}{2}x-\dfrac{1}{2})$ với $x\in\mathbb R$.}
\end{cauhoi}
\begin{cauhoi}{Bài 23}{9H1-C1-B1-D02-TL-TH-023}
Viết nghiệm tổng quát của phương trình $2x-3y=5$ (biểu diễn $y$ theo $x$).
\loigiai{Ta có $-3y=5-2x$ nên $y=\dfrac{2}{3}x-\dfrac{5}{3}$. Nghiệm tổng quát là $(x;\,\dfrac{2}{3}x-\dfrac{5}{3})$ với $x\in\mathbb R$.}
\end{cauhoi}
\begin{cauhoi}{Bài 24}{9H1-C1-B1-D02-TL-TH-024}
Viết nghiệm tổng quát của phương trình $-x+3y=-3$ (biểu diễn $y$ theo $x$).
\loigiai{Ta có $3y=-3+x$ nên $y=\dfrac{1}{3}x-1$. Nghiệm tổng quát là $(x;\,\dfrac{1}{3}x-1)$ với $x\in\mathbb R$.}
\end{cauhoi}
\begin{cauhoi}{Bài 25}{9H1-C1-B1-D02-TL-TH-025}
Viết nghiệm tổng quát của phương trình $5x+2y=8$ (biểu diễn $y$ theo $x$).
\loigiai{Ta có $2y=8-5x$ nên $y=-\dfrac{5}{2}x+4$. Nghiệm tổng quát là $(x;\,-\dfrac{5}{2}x+4)$ với $x\in\mathbb R$.}
\end{cauhoi}
\begin{cauhoi}{Bài 26}{9H1-C1-B1-D02-TL-TH-026}
Viết nghiệm tổng quát của phương trình $-5x-3y=2$ (biểu diễn $y$ theo $x$).
\loigiai{Ta có $-3y=2+5x$ nên $y=-\dfrac{5}{3}x-\dfrac{2}{3}$. Nghiệm tổng quát là $(x;\,-\dfrac{5}{3}x-\dfrac{2}{3})$ với $x\in\mathbb R$.}
\end{cauhoi}
\begin{cauhoi}{Bài 27}{9H1-C1-B1-D02-TL-TH-027}
Viết nghiệm tổng quát của phương trình $x+4y=8$ (biểu diễn $y$ theo $x$).
\loigiai{Ta có $4y=8-x$ nên $y=-\dfrac{1}{4}x+2$. Nghiệm tổng quát là $(x;\,-\dfrac{1}{4}x+2)$ với $x\in\mathbb R$.}
\end{cauhoi}
\begin{cauhoi}{Bài 28}{9H1-C1-B1-D02-TL-TH-028}
Viết nghiệm tổng quát của phương trình $5x-3y=7$ (biểu diễn $y$ theo $x$).
\loigiai{Ta có $-3y=7-5x$ nên $y=\dfrac{5}{3}x-\dfrac{7}{3}$. Nghiệm tổng quát là $(x;\,\dfrac{5}{3}x-\dfrac{7}{3})$ với $x\in\mathbb R$.}
\end{cauhoi}
\begin{cauhoi}{Bài 29}{9H1-C1-B1-D02-TL-TH-029}
Viết nghiệm tổng quát của phương trình $-x+3y=-6$ (biểu diễn $y$ theo $x$).
\loigiai{Ta có $3y=-6+x$ nên $y=\dfrac{1}{3}x-2$. Nghiệm tổng quát là $(x;\,\dfrac{1}{3}x-2)$ với $x\in\mathbb R$.}
\end{cauhoi}
\begin{cauhoi}{Bài 30}{9H1-C1-B1-D02-TL-TH-030}
Viết nghiệm tổng quát của phương trình $-2x+3y=-4$ (biểu diễn $y$ theo $x$).
\loigiai{Ta có $3y=-4+2x$ nên $y=\dfrac{2}{3}x-\dfrac{4}{3}$. Nghiệm tổng quát là $(x;\,\dfrac{2}{3}x-\dfrac{4}{3})$ với $x\in\mathbb R$.}
\end{cauhoi}
\begin{cauhoi}{Bài 31}{9H1-C1-B1-D02-TL-TH-031}
Viết nghiệm tổng quát của phương trình $-5x+4y=3$ (biểu diễn $y$ theo $x$).
\loigiai{Ta có $4y=3+5x$ nên $y=\dfrac{5}{4}x+\dfrac{3}{4}$. Nghiệm tổng quát là $(x;\,\dfrac{5}{4}x+\dfrac{3}{4})$ với $x\in\mathbb R$.}
\end{cauhoi}
\begin{cauhoi}{Bài 32}{9H1-C1-B1-D02-TL-TH-032}
Viết nghiệm tổng quát của phương trình $-2x+4y=-3$ (biểu diễn $y$ theo $x$).
\loigiai{Ta có $4y=-3+2x$ nên $y=\dfrac{1}{2}x-\dfrac{3}{4}$. Nghiệm tổng quát là $(x;\,\dfrac{1}{2}x-\dfrac{3}{4})$ với $x\in\mathbb R$.}
\end{cauhoi}
\begin{cauhoi}{Bài 33}{9H1-C1-B1-D02-TL-TH-033}
Viết nghiệm tổng quát của phương trình $-x-3y=-7$ (biểu diễn $y$ theo $x$).
\loigiai{Ta có $-3y=-7+x$ nên $y=-\dfrac{1}{3}x+\dfrac{7}{3}$. Nghiệm tổng quát là $(x;\,-\dfrac{1}{3}x+\dfrac{7}{3})$ với $x\in\mathbb R$.}
\end{cauhoi}
\begin{cauhoi}{Bài 34}{9H1-C1-B1-D02-TL-TH-034}
Viết nghiệm tổng quát của phương trình $x+5y=8$ (biểu diễn $y$ theo $x$).
\loigiai{Ta có $5y=8-x$ nên $y=-\dfrac{1}{5}x+\dfrac{8}{5}$. Nghiệm tổng quát là $(x;\,-\dfrac{1}{5}x+\dfrac{8}{5})$ với $x\in\mathbb R$.}
\end{cauhoi}
\begin{cauhoi}{Bài 35}{9H1-C1-B1-D02-TL-TH-035}
Viết nghiệm tổng quát của phương trình $x+3y=-7$ (biểu diễn $y$ theo $x$).
\loigiai{Ta có $3y=-7-x$ nên $y=-\dfrac{1}{3}x-\dfrac{7}{3}$. Nghiệm tổng quát là $(x;\,-\dfrac{1}{3}x-\dfrac{7}{3})$ với $x\in\mathbb R$.}
\end{cauhoi}
\begin{cauhoi}{Bài 36}{9H1-C1-B1-D02-TL-TH-036}
Viết nghiệm tổng quát của phương trình $-x-2y=-8$ (biểu diễn $y$ theo $x$).
\loigiai{Ta có $-2y=-8+x$ nên $y=-\dfrac{1}{2}x+4$. Nghiệm tổng quát là $(x;\,-\dfrac{1}{2}x+4)$ với $x\in\mathbb R$.}
\end{cauhoi}
\begin{cauhoi}{Bài 37}{9H1-C1-B1-D02-TL-TH-037}
Viết nghiệm tổng quát của phương trình $-5x-3y=-8$ (biểu diễn $y$ theo $x$).
\loigiai{Ta có $-3y=-8+5x$ nên $y=-\dfrac{5}{3}x+\dfrac{8}{3}$. Nghiệm tổng quát là $(x;\,-\dfrac{5}{3}x+\dfrac{8}{3})$ với $x\in\mathbb R$.}
\end{cauhoi}
\begin{cauhoi}{Bài 38}{9H1-C1-B1-D02-TL-TH-038}
Viết nghiệm tổng quát của phương trình $2x-3y=-2$ (biểu diễn $y$ theo $x$).
\loigiai{Ta có $-3y=-2-2x$ nên $y=\dfrac{2}{3}x+\dfrac{2}{3}$. Nghiệm tổng quát là $(x;\,\dfrac{2}{3}x+\dfrac{2}{3})$ với $x\in\mathbb R$.}
\end{cauhoi}
\begin{cauhoi}{Bài 39}{9H1-C1-B1-D02-TL-TH-039}
Viết nghiệm tổng quát của phương trình $-4x+5y=-6$ (biểu diễn $y$ theo $x$).
\loigiai{Ta có $5y=-6+4x$ nên $y=\dfrac{4}{5}x-\dfrac{6}{5}$. Nghiệm tổng quát là $(x;\,\dfrac{4}{5}x-\dfrac{6}{5})$ với $x\in\mathbb R$.}
\end{cauhoi}
\begin{cauhoi}{Bài 40}{9H1-C1-B1-D02-TL-TH-040}
Viết nghiệm tổng quát của phương trình $2x-3y=-7$ (biểu diễn $y$ theo $x$).
\loigiai{Ta có $-3y=-7-2x$ nên $y=\dfrac{2}{3}x+\dfrac{7}{3}$. Nghiệm tổng quát là $(x;\,\dfrac{2}{3}x+\dfrac{7}{3})$ với $x\in\mathbb R$.}
\end{cauhoi}
\begin{cauhoi}{Bài 41}{9H1-C1-B1-D02-TL-TH-041}
Viết nghiệm tổng quát của phương trình $-2x+5y=-2$ (biểu diễn $y$ theo $x$).
\loigiai{Ta có $5y=-2+2x$ nên $y=\dfrac{2}{5}x-\dfrac{2}{5}$. Nghiệm tổng quát là $(x;\,\dfrac{2}{5}x-\dfrac{2}{5})$ với $x\in\mathbb R$.}
\end{cauhoi}
\begin{cauhoi}{Bài 42}{9H1-C1-B1-D02-TL-TH-042}
Viết nghiệm tổng quát của phương trình $-3x+5y=-3$ (biểu diễn $y$ theo $x$).
\loigiai{Ta có $5y=-3+3x$ nên $y=\dfrac{3}{5}x-\dfrac{3}{5}$. Nghiệm tổng quát là $(x;\,\dfrac{3}{5}x-\dfrac{3}{5})$ với $x\in\mathbb R$.}
\end{cauhoi}
\begin{cauhoi}{Bài 43}{9H1-C1-B1-D02-TL-TH-043}
Viết nghiệm tổng quát của phương trình $-3x-2y=6$ (biểu diễn $y$ theo $x$).
\loigiai{Ta có $-2y=6+3x$ nên $y=-\dfrac{3}{2}x-3$. Nghiệm tổng quát là $(x;\,-\dfrac{3}{2}x-3)$ với $x\in\mathbb R$.}
\end{cauhoi}
\begin{cauhoi}{Bài 44}{9H1-C1-B1-D02-TL-TH-044}
Viết nghiệm tổng quát của phương trình $-x-2y=-2$ (biểu diễn $y$ theo $x$).
\loigiai{Ta có $-2y=-2+x$ nên $y=-\dfrac{1}{2}x+1$. Nghiệm tổng quát là $(x;\,-\dfrac{1}{2}x+1)$ với $x\in\mathbb R$.}
\end{cauhoi}
\begin{cauhoi}{Bài 45}{9H1-C1-B1-D02-TL-TH-045}
Viết nghiệm tổng quát của phương trình $5x-2y=-9$ (biểu diễn $y$ theo $x$).
\loigiai{Ta có $-2y=-9-5x$ nên $y=\dfrac{5}{2}x+\dfrac{9}{2}$. Nghiệm tổng quát là $(x;\,\dfrac{5}{2}x+\dfrac{9}{2})$ với $x\in\mathbb R$.}
\end{cauhoi}
\begin{cauhoi}{Bài 46}{9H1-C1-B1-D02-TL-TH-046}
Viết nghiệm tổng quát của phương trình $-x+2y=2$ (biểu diễn $y$ theo $x$).
\loigiai{Ta có $2y=2+x$ nên $y=\dfrac{1}{2}x+1$. Nghiệm tổng quát là $(x;\,\dfrac{1}{2}x+1)$ với $x\in\mathbb R$.}
\end{cauhoi}
\begin{cauhoi}{Bài 47}{9H1-C1-B1-D02-TL-TH-047}
Viết nghiệm tổng quát của phương trình $-3x+5y=-7$ (biểu diễn $y$ theo $x$).
\loigiai{Ta có $5y=-7+3x$ nên $y=\dfrac{3}{5}x-\dfrac{7}{5}$. Nghiệm tổng quát là $(x;\,\dfrac{3}{5}x-\dfrac{7}{5})$ với $x\in\mathbb R$.}
\end{cauhoi}
\begin{cauhoi}{Bài 48}{9H1-C1-B1-D02-TL-TH-048}
Viết nghiệm tổng quát của phương trình $3x+2y=5$ (biểu diễn $y$ theo $x$).
\loigiai{Ta có $2y=5-3x$ nên $y=-\dfrac{3}{2}x+\dfrac{5}{2}$. Nghiệm tổng quát là $(x;\,-\dfrac{3}{2}x+\dfrac{5}{2})$ với $x\in\mathbb R$.}
\end{cauhoi}
\begin{cauhoi}{Bài 49}{9H1-C1-B1-D02-TL-TH-049}
Viết nghiệm tổng quát của phương trình $5x+4y=3$ (biểu diễn $y$ theo $x$).
\loigiai{Ta có $4y=3-5x$ nên $y=-\dfrac{5}{4}x+\dfrac{3}{4}$. Nghiệm tổng quát là $(x;\,-\dfrac{5}{4}x+\dfrac{3}{4})$ với $x\in\mathbb R$.}
\end{cauhoi}
\begin{cauhoi}{Bài 50}{9H1-C1-B1-D02-TL-TH-050}
Viết nghiệm tổng quát của phương trình $2x+4y=2$ (biểu diễn $y$ theo $x$).
\loigiai{Ta có $4y=2-2x$ nên $y=-\dfrac{1}{2}x+\dfrac{1}{2}$. Nghiệm tổng quát là $(x;\,-\dfrac{1}{2}x+\dfrac{1}{2})$ với $x\in\mathbb R$.}
\end{cauhoi}
\clearpage\section*{Phần II. Trắc nghiệm}
\begin{cauhoi}{Câu 1}{9H1-C1-B1-D02-TN-TH-001}
Viết nghiệm tổng quát của phương trình $-x+5y=-9$ (biểu diễn $y$ theo $x$).
\begin{tasks}(2)
\task $y=5x+9$
\task $y=-\dfrac{1}{5}x-\dfrac{9}{5}$
\task $y=\dfrac{1}{5}x-\dfrac{9}{5}$
\task $y=\dfrac{1}{5}x+\dfrac{9}{5}$
\end{tasks}
\dapan{C}
\end{cauhoi}
\begin{cauhoi}{Câu 2}{9H1-C1-B1-D02-TN-TH-002}
Viết nghiệm tổng quát của phương trình $-4x+3y=-5$ (biểu diễn $y$ theo $x$).
\begin{tasks}(2)
\task $y=\dfrac{4}{3}x-\dfrac{5}{3}$
\task $y=-\dfrac{4}{3}x-\dfrac{5}{3}$
\task $y=\dfrac{3}{4}x+\dfrac{5}{4}$
\task $y=\dfrac{4}{3}x+\dfrac{5}{3}$
\end{tasks}
\dapan{A}
\end{cauhoi}
\begin{cauhoi}{Câu 3}{9H1-C1-B1-D02-TN-TH-003}
Viết nghiệm tổng quát của phương trình $2x-3y=6$ (biểu diễn $y$ theo $x$).
\begin{tasks}(2)
\task $y=\dfrac{3}{2}x+3$
\task $y=\dfrac{2}{3}x-2$
\task $y=\dfrac{2}{3}x+2$
\task $y=-\dfrac{2}{3}x-2$
\end{tasks}
\dapan{B}
\end{cauhoi}
\begin{cauhoi}{Câu 4}{9H1-C1-B1-D02-TN-TH-004}
Viết nghiệm tổng quát của phương trình $-5x+2y=-1$ (biểu diễn $y$ theo $x$).
\begin{tasks}(2)
\task $y=\dfrac{2}{5}x+\dfrac{1}{5}$
\task $y=\dfrac{5}{2}x+\dfrac{1}{2}$
\task $y=\dfrac{5}{2}x-\dfrac{1}{2}$
\task $y=-\dfrac{5}{2}x-\dfrac{1}{2}$
\end{tasks}
\dapan{C}
\end{cauhoi}
\begin{cauhoi}{Câu 5}{9H1-C1-B1-D02-TN-TH-005}
Viết nghiệm tổng quát của phương trình $3x+5y=3$ (biểu diễn $y$ theo $x$).
\begin{tasks}(2)
\task $y=-\dfrac{5}{3}x+1$
\task $y=-\dfrac{3}{5}x+\dfrac{3}{5}$
\task $y=\dfrac{3}{5}x+\dfrac{3}{5}$
\task $y=-\dfrac{3}{5}x-\dfrac{3}{5}$
\end{tasks}
\dapan{B}
\end{cauhoi}
\begin{cauhoi}{Câu 6}{9H1-C1-B1-D02-TN-TH-006}
Viết nghiệm tổng quát của phương trình $4x+3y=7$ (biểu diễn $y$ theo $x$).
\begin{tasks}(2)
\task $y=-\dfrac{4}{3}x-\dfrac{7}{3}$
\task $y=-\dfrac{3}{4}x+\dfrac{7}{4}$
\task $y=-\dfrac{4}{3}x+\dfrac{7}{3}$
\task $y=\dfrac{4}{3}x+\dfrac{7}{3}$
\end{tasks}
\dapan{C}
\end{cauhoi}
\begin{cauhoi}{Câu 7}{9H1-C1-B1-D02-TN-TH-007}
Viết nghiệm tổng quát của phương trình $-2x-3y=-4$ (biểu diễn $y$ theo $x$).
\begin{tasks}(2)
\task $y=-\dfrac{3}{2}x+2$
\task $y=\dfrac{2}{3}x+\dfrac{4}{3}$
\task $y=-\dfrac{2}{3}x-\dfrac{4}{3}$
\task $y=-\dfrac{2}{3}x+\dfrac{4}{3}$
\end{tasks}
\dapan{D}
\end{cauhoi}
\begin{cauhoi}{Câu 8}{9H1-C1-B1-D02-TN-TH-008}
Viết nghiệm tổng quát của phương trình $-3x+5y=4$ (biểu diễn $y$ theo $x$).
\begin{tasks}(2)
\task $y=\dfrac{3}{5}x-\dfrac{4}{5}$
\task $y=\dfrac{3}{5}x+\dfrac{4}{5}$
\task $y=-\dfrac{3}{5}x+\dfrac{4}{5}$
\task $y=\dfrac{5}{3}x-\dfrac{4}{3}$
\end{tasks}
\dapan{B}
\end{cauhoi}
\begin{cauhoi}{Câu 9}{9H1-C1-B1-D02-TN-TH-009}
Viết nghiệm tổng quát của phương trình $2x+5y=-5$ (biểu diễn $y$ theo $x$).
\begin{tasks}(2)
\task $y=\dfrac{2}{5}x-1$
\task $y=-\dfrac{5}{2}x-\dfrac{5}{2}$
\task $y=-\dfrac{2}{5}x-1$
\task $y=-\dfrac{2}{5}x+1$
\end{tasks}
\dapan{C}
\end{cauhoi}
\begin{cauhoi}{Câu 10}{9H1-C1-B1-D02-TN-TH-010}
Viết nghiệm tổng quát của phương trình $2x+5y=-6$ (biểu diễn $y$ theo $x$).
\begin{tasks}(2)
\task $y=-\dfrac{5}{2}x-3$
\task $y=-\dfrac{2}{5}x-\dfrac{6}{5}$
\task $y=-\dfrac{2}{5}x+\dfrac{6}{5}$
\task $y=\dfrac{2}{5}x-\dfrac{6}{5}$
\end{tasks}
\dapan{B}
\end{cauhoi}
\begin{cauhoi}{Câu 11}{9H1-C1-B1-D02-TN-TH-011}
Viết nghiệm tổng quát của phương trình $-3x+4y=5$ (biểu diễn $y$ theo $x$).
\begin{tasks}(2)
\task $y=\dfrac{3}{4}x-\dfrac{5}{4}$
\task $y=-\dfrac{3}{4}x+\dfrac{5}{4}$
\task $y=\dfrac{4}{3}x-\dfrac{5}{3}$
\task $y=\dfrac{3}{4}x+\dfrac{5}{4}$
\end{tasks}
\dapan{D}
\end{cauhoi}
\begin{cauhoi}{Câu 12}{9H1-C1-B1-D02-TN-TH-012}
Viết nghiệm tổng quát của phương trình $4x+5y=5$ (biểu diễn $y$ theo $x$).
\begin{tasks}(2)
\task $y=-\dfrac{4}{5}x-1$
\task $y=-\dfrac{4}{5}x+1$
\task $y=\dfrac{4}{5}x+1$
\task $y=-\dfrac{5}{4}x+\dfrac{5}{4}$
\end{tasks}
\dapan{B}
\end{cauhoi}
\begin{cauhoi}{Câu 13}{9H1-C1-B1-D02-TN-TH-013}
Viết nghiệm tổng quát của phương trình $2x-3y=-1$ (biểu diễn $y$ theo $x$).
\begin{tasks}(2)
\task $y=-\dfrac{2}{3}x+\dfrac{1}{3}$
\task $y=\dfrac{2}{3}x+\dfrac{1}{3}$
\task $y=\dfrac{2}{3}x-\dfrac{1}{3}$
\task $y=\dfrac{3}{2}x-\dfrac{1}{2}$
\end{tasks}
\dapan{B}
\end{cauhoi}
\begin{cauhoi}{Câu 14}{9H1-C1-B1-D02-TN-TH-014}
Viết nghiệm tổng quát của phương trình $x+2y=4$ (biểu diễn $y$ theo $x$).
\begin{tasks}(2)
\task $y=-2x+4$
\task $y=-\dfrac{1}{2}x+2$
\task $y=\dfrac{1}{2}x+2$
\task $y=-\dfrac{1}{2}x-2$
\end{tasks}
\dapan{B}
\end{cauhoi}
\begin{cauhoi}{Câu 15}{9H1-C1-B1-D02-TN-TH-015}
Viết nghiệm tổng quát của phương trình $-x+4y=-3$ (biểu diễn $y$ theo $x$).
\begin{tasks}(2)
\task $y=4x+3$
\task $y=\dfrac{1}{4}x-\dfrac{3}{4}$
\task $y=\dfrac{1}{4}x+\dfrac{3}{4}$
\task $y=-\dfrac{1}{4}x-\dfrac{3}{4}$
\end{tasks}
\dapan{B}
\end{cauhoi}
\begin{cauhoi}{Câu 16}{9H1-C1-B1-D02-TN-TH-016}
Viết nghiệm tổng quát của phương trình $2x-3y=8$ (biểu diễn $y$ theo $x$).
\begin{tasks}(2)
\task $y=\dfrac{2}{3}x-\dfrac{8}{3}$
\task $y=-\dfrac{2}{3}x-\dfrac{8}{3}$
\task $y=\dfrac{3}{2}x+4$
\task $y=\dfrac{2}{3}x+\dfrac{8}{3}$
\end{tasks}
\dapan{A}
\end{cauhoi}
\begin{cauhoi}{Câu 17}{9H1-C1-B1-D02-TN-TH-017}
Viết nghiệm tổng quát của phương trình $-x+5y=-1$ (biểu diễn $y$ theo $x$).
\begin{tasks}(2)
\task $y=\dfrac{1}{5}x+\dfrac{1}{5}$
\task $y=-\dfrac{1}{5}x-\dfrac{1}{5}$
\task $y=\dfrac{1}{5}x-\dfrac{1}{5}$
\task $y=5x+1$
\end{tasks}
\dapan{C}
\end{cauhoi}
\begin{cauhoi}{Câu 18}{9H1-C1-B1-D02-TN-TH-018}
Viết nghiệm tổng quát của phương trình $5x+4y=8$ (biểu diễn $y$ theo $x$).
\begin{tasks}(2)
\task $y=-\dfrac{5}{4}x-2$
\task $y=\dfrac{5}{4}x+2$
\task $y=-\dfrac{4}{5}x+\dfrac{8}{5}$
\task $y=-\dfrac{5}{4}x+2$
\end{tasks}
\dapan{D}
\end{cauhoi}
\begin{cauhoi}{Câu 19}{9H1-C1-B1-D02-TN-TH-019}
Viết nghiệm tổng quát của phương trình $-x+2y=5$ (biểu diễn $y$ theo $x$).
\begin{tasks}(2)
\task $y=\dfrac{1}{2}x-\dfrac{5}{2}$
\task $y=-\dfrac{1}{2}x+\dfrac{5}{2}$
\task $y=2x-5$
\task $y=\dfrac{1}{2}x+\dfrac{5}{2}$
\end{tasks}
\dapan{D}
\end{cauhoi}
\begin{cauhoi}{Câu 20}{9H1-C1-B1-D02-TN-TH-020}
Viết nghiệm tổng quát của phương trình $-2x+4y=-7$ (biểu diễn $y$ theo $x$).
\begin{tasks}(2)
\task $y=\dfrac{1}{2}x-\dfrac{7}{4}$
\task $y=2x+\dfrac{7}{2}$
\task $y=\dfrac{1}{2}x+\dfrac{7}{4}$
\task $y=-\dfrac{1}{2}x-\dfrac{7}{4}$
\end{tasks}
\dapan{A}
\end{cauhoi}
\begin{cauhoi}{Câu 21}{9H1-C1-B1-D02-TN-TH-021}
Viết nghiệm tổng quát của phương trình $4x-2y=1$ (biểu diễn $y$ theo $x$).
\begin{tasks}(2)
\task $y=\dfrac{1}{2}x+\dfrac{1}{4}$
\task $y=2x+\dfrac{1}{2}$
\task $y=2x-\dfrac{1}{2}$
\task $y=-2x-\dfrac{1}{2}$
\end{tasks}
\dapan{C}
\end{cauhoi}
\begin{cauhoi}{Câu 22}{9H1-C1-B1-D02-TN-TH-022}
Viết nghiệm tổng quát của phương trình $3x+2y=-1$ (biểu diễn $y$ theo $x$).
\begin{tasks}(2)
\task $y=-\dfrac{2}{3}x-\dfrac{1}{3}$
\task $y=\dfrac{3}{2}x-\dfrac{1}{2}$
\task $y=-\dfrac{3}{2}x+\dfrac{1}{2}$
\task $y=-\dfrac{3}{2}x-\dfrac{1}{2}$
\end{tasks}
\dapan{D}
\end{cauhoi}
\begin{cauhoi}{Câu 23}{9H1-C1-B1-D02-TN-TH-023}
Viết nghiệm tổng quát của phương trình $2x-3y=5$ (biểu diễn $y$ theo $x$).
\begin{tasks}(2)
\task $y=\dfrac{3}{2}x+\dfrac{5}{2}$
\task $y=-\dfrac{2}{3}x-\dfrac{5}{3}$
\task $y=\dfrac{2}{3}x+\dfrac{5}{3}$
\task $y=\dfrac{2}{3}x-\dfrac{5}{3}$
\end{tasks}
\dapan{D}
\end{cauhoi}
\begin{cauhoi}{Câu 24}{9H1-C1-B1-D02-TN-TH-024}
Viết nghiệm tổng quát của phương trình $-x+3y=-3$ (biểu diễn $y$ theo $x$).
\begin{tasks}(2)
\task $y=\dfrac{1}{3}x-1$
\task $y=\dfrac{1}{3}x+1$
\task $y=3x+3$
\task $y=-\dfrac{1}{3}x-1$
\end{tasks}
\dapan{A}
\end{cauhoi}
\begin{cauhoi}{Câu 25}{9H1-C1-B1-D02-TN-TH-025}
Viết nghiệm tổng quát của phương trình $5x+2y=8$ (biểu diễn $y$ theo $x$).
\begin{tasks}(2)
\task $y=\dfrac{5}{2}x+4$
\task $y=-\dfrac{5}{2}x+4$
\task $y=-\dfrac{5}{2}x-4$
\task $y=-\dfrac{2}{5}x+\dfrac{8}{5}$
\end{tasks}
\dapan{B}
\end{cauhoi}
\begin{cauhoi}{Câu 26}{9H1-C1-B1-D02-TN-TH-026}
Viết nghiệm tổng quát của phương trình $-5x-3y=2$ (biểu diễn $y$ theo $x$).
\begin{tasks}(2)
\task $y=\dfrac{5}{3}x-\dfrac{2}{3}$
\task $y=-\dfrac{5}{3}x+\dfrac{2}{3}$
\task $y=-\dfrac{3}{5}x-\dfrac{2}{5}$
\task $y=-\dfrac{5}{3}x-\dfrac{2}{3}$
\end{tasks}
\dapan{D}
\end{cauhoi}
\begin{cauhoi}{Câu 27}{9H1-C1-B1-D02-TN-TH-027}
Viết nghiệm tổng quát của phương trình $x+4y=8$ (biểu diễn $y$ theo $x$).
\begin{tasks}(2)
\task $y=-4x+8$
\task $y=-\dfrac{1}{4}x-2$
\task $y=-\dfrac{1}{4}x+2$
\task $y=\dfrac{1}{4}x+2$
\end{tasks}
\dapan{C}
\end{cauhoi}
\begin{cauhoi}{Câu 28}{9H1-C1-B1-D02-TN-TH-028}
Viết nghiệm tổng quát của phương trình $5x-3y=7$ (biểu diễn $y$ theo $x$).
\begin{tasks}(2)
\task $y=-\dfrac{5}{3}x-\dfrac{7}{3}$
\task $y=\dfrac{3}{5}x+\dfrac{7}{5}$
\task $y=\dfrac{5}{3}x-\dfrac{7}{3}$
\task $y=\dfrac{5}{3}x+\dfrac{7}{3}$
\end{tasks}
\dapan{C}
\end{cauhoi}
\begin{cauhoi}{Câu 29}{9H1-C1-B1-D02-TN-TH-029}
Viết nghiệm tổng quát của phương trình $-x+3y=-6$ (biểu diễn $y$ theo $x$).
\begin{tasks}(2)
\task $y=3x+6$
\task $y=-\dfrac{1}{3}x-2$
\task $y=\dfrac{1}{3}x-2$
\task $y=\dfrac{1}{3}x+2$
\end{tasks}
\dapan{C}
\end{cauhoi}
\begin{cauhoi}{Câu 30}{9H1-C1-B1-D02-TN-TH-030}
Viết nghiệm tổng quát của phương trình $-2x+3y=-4$ (biểu diễn $y$ theo $x$).
\begin{tasks}(2)
\task $y=\dfrac{3}{2}x+2$
\task $y=\dfrac{2}{3}x-\dfrac{4}{3}$
\task $y=-\dfrac{2}{3}x-\dfrac{4}{3}$
\task $y=\dfrac{2}{3}x+\dfrac{4}{3}$
\end{tasks}
\dapan{B}
\end{cauhoi}
\begin{cauhoi}{Câu 31}{9H1-C1-B1-D02-TN-TH-031}
Viết nghiệm tổng quát của phương trình $-5x+4y=3$ (biểu diễn $y$ theo $x$).
\begin{tasks}(2)
\task $y=\dfrac{5}{4}x-\dfrac{3}{4}$
\task $y=\dfrac{5}{4}x+\dfrac{3}{4}$
\task $y=\dfrac{4}{5}x-\dfrac{3}{5}$
\task $y=-\dfrac{5}{4}x+\dfrac{3}{4}$
\end{tasks}
\dapan{B}
\end{cauhoi}
\begin{cauhoi}{Câu 32}{9H1-C1-B1-D02-TN-TH-032}
Viết nghiệm tổng quát của phương trình $-2x+4y=-3$ (biểu diễn $y$ theo $x$).
\begin{tasks}(2)
\task $y=\dfrac{1}{2}x+\dfrac{3}{4}$
\task $y=2x+\dfrac{3}{2}$
\task $y=-\dfrac{1}{2}x-\dfrac{3}{4}$
\task $y=\dfrac{1}{2}x-\dfrac{3}{4}$
\end{tasks}
\dapan{D}
\end{cauhoi}
\begin{cauhoi}{Câu 33}{9H1-C1-B1-D02-TN-TH-033}
Viết nghiệm tổng quát của phương trình $-x-3y=-7$ (biểu diễn $y$ theo $x$).
\begin{tasks}(2)
\task $y=-\dfrac{1}{3}x-\dfrac{7}{3}$
\task $y=\dfrac{1}{3}x+\dfrac{7}{3}$
\task $y=-3x+7$
\task $y=-\dfrac{1}{3}x+\dfrac{7}{3}$
\end{tasks}
\dapan{D}
\end{cauhoi}
\begin{cauhoi}{Câu 34}{9H1-C1-B1-D02-TN-TH-034}
Viết nghiệm tổng quát của phương trình $x+5y=8$ (biểu diễn $y$ theo $x$).
\begin{tasks}(2)
\task $y=-5x+8$
\task $y=-\dfrac{1}{5}x-\dfrac{8}{5}$
\task $y=\dfrac{1}{5}x+\dfrac{8}{5}$
\task $y=-\dfrac{1}{5}x+\dfrac{8}{5}$
\end{tasks}
\dapan{D}
\end{cauhoi}
\begin{cauhoi}{Câu 35}{9H1-C1-B1-D02-TN-TH-035}
Viết nghiệm tổng quát của phương trình $x+3y=-7$ (biểu diễn $y$ theo $x$).
\begin{tasks}(2)
\task $y=-3x-7$
\task $y=-\dfrac{1}{3}x-\dfrac{7}{3}$
\task $y=-\dfrac{1}{3}x+\dfrac{7}{3}$
\task $y=\dfrac{1}{3}x-\dfrac{7}{3}$
\end{tasks}
\dapan{B}
\end{cauhoi}
\begin{cauhoi}{Câu 36}{9H1-C1-B1-D02-TN-TH-036}
Viết nghiệm tổng quát của phương trình $-x-2y=-8$ (biểu diễn $y$ theo $x$).
\begin{tasks}(2)
\task $y=-\dfrac{1}{2}x-4$
\task $y=-\dfrac{1}{2}x+4$
\task $y=-2x+8$
\task $y=\dfrac{1}{2}x+4$
\end{tasks}
\dapan{B}
\end{cauhoi}
\begin{cauhoi}{Câu 37}{9H1-C1-B1-D02-TN-TH-037}
Viết nghiệm tổng quát của phương trình $-5x-3y=-8$ (biểu diễn $y$ theo $x$).
\begin{tasks}(2)
\task $y=\dfrac{5}{3}x+\dfrac{8}{3}$
\task $y=-\dfrac{3}{5}x+\dfrac{8}{5}$
\task $y=-\dfrac{5}{3}x+\dfrac{8}{3}$
\task $y=-\dfrac{5}{3}x-\dfrac{8}{3}$
\end{tasks}
\dapan{C}
\end{cauhoi}
\begin{cauhoi}{Câu 38}{9H1-C1-B1-D02-TN-TH-038}
Viết nghiệm tổng quát của phương trình $2x-3y=-2$ (biểu diễn $y$ theo $x$).
\begin{tasks}(2)
\task $y=\dfrac{2}{3}x-\dfrac{2}{3}$
\task $y=\dfrac{3}{2}x-1$
\task $y=\dfrac{2}{3}x+\dfrac{2}{3}$
\task $y=-\dfrac{2}{3}x+\dfrac{2}{3}$
\end{tasks}
\dapan{C}
\end{cauhoi}
\begin{cauhoi}{Câu 39}{9H1-C1-B1-D02-TN-TH-039}
Viết nghiệm tổng quát của phương trình $-4x+5y=-6$ (biểu diễn $y$ theo $x$).
\begin{tasks}(2)
\task $y=\dfrac{4}{5}x+\dfrac{6}{5}$
\task $y=\dfrac{4}{5}x-\dfrac{6}{5}$
\task $y=\dfrac{5}{4}x+\dfrac{3}{2}$
\task $y=-\dfrac{4}{5}x-\dfrac{6}{5}$
\end{tasks}
\dapan{B}
\end{cauhoi}
\begin{cauhoi}{Câu 40}{9H1-C1-B1-D02-TN-TH-040}
Viết nghiệm tổng quát của phương trình $2x-3y=-7$ (biểu diễn $y$ theo $x$).
\begin{tasks}(2)
\task $y=\dfrac{3}{2}x-\dfrac{7}{2}$
\task $y=-\dfrac{2}{3}x+\dfrac{7}{3}$
\task $y=\dfrac{2}{3}x+\dfrac{7}{3}$
\task $y=\dfrac{2}{3}x-\dfrac{7}{3}$
\end{tasks}
\dapan{C}
\end{cauhoi}
\begin{cauhoi}{Câu 41}{9H1-C1-B1-D02-TN-TH-041}
Viết nghiệm tổng quát của phương trình $-2x+5y=-2$ (biểu diễn $y$ theo $x$).
\begin{tasks}(2)
\task $y=\dfrac{5}{2}x+1$
\task $y=\dfrac{2}{5}x-\dfrac{2}{5}$
\task $y=\dfrac{2}{5}x+\dfrac{2}{5}$
\task $y=-\dfrac{2}{5}x-\dfrac{2}{5}$
\end{tasks}
\dapan{B}
\end{cauhoi}
\begin{cauhoi}{Câu 42}{9H1-C1-B1-D02-TN-TH-042}
Viết nghiệm tổng quát của phương trình $-3x+5y=-3$ (biểu diễn $y$ theo $x$).
\begin{tasks}(2)
\task $y=\dfrac{3}{5}x+\dfrac{3}{5}$
\task $y=\dfrac{3}{5}x-\dfrac{3}{5}$
\task $y=-\dfrac{3}{5}x-\dfrac{3}{5}$
\task $y=\dfrac{5}{3}x+1$
\end{tasks}
\dapan{B}
\end{cauhoi}
\begin{cauhoi}{Câu 43}{9H1-C1-B1-D02-TN-TH-043}
Viết nghiệm tổng quát của phương trình $-3x-2y=6$ (biểu diễn $y$ theo $x$).
\begin{tasks}(2)
\task $y=-\dfrac{3}{2}x-3$
\task $y=\dfrac{3}{2}x-3$
\task $y=-\dfrac{2}{3}x-2$
\task $y=-\dfrac{3}{2}x+3$
\end{tasks}
\dapan{A}
\end{cauhoi}
\begin{cauhoi}{Câu 44}{9H1-C1-B1-D02-TN-TH-044}
Viết nghiệm tổng quát của phương trình $-x-2y=-2$ (biểu diễn $y$ theo $x$).
\begin{tasks}(2)
\task $y=-2x+2$
\task $y=-\dfrac{1}{2}x-1$
\task $y=-\dfrac{1}{2}x+1$
\task $y=\dfrac{1}{2}x+1$
\end{tasks}
\dapan{C}
\end{cauhoi}
\begin{cauhoi}{Câu 45}{9H1-C1-B1-D02-TN-TH-045}
Viết nghiệm tổng quát của phương trình $5x-2y=-9$ (biểu diễn $y$ theo $x$).
\begin{tasks}(2)
\task $y=\dfrac{2}{5}x-\dfrac{9}{5}$
\task $y=-\dfrac{5}{2}x+\dfrac{9}{2}$
\task $y=\dfrac{5}{2}x-\dfrac{9}{2}$
\task $y=\dfrac{5}{2}x+\dfrac{9}{2}$
\end{tasks}
\dapan{D}
\end{cauhoi}
\begin{cauhoi}{Câu 46}{9H1-C1-B1-D02-TN-TH-046}
Viết nghiệm tổng quát của phương trình $-x+2y=2$ (biểu diễn $y$ theo $x$).
\begin{tasks}(2)
\task $y=\dfrac{1}{2}x+1$
\task $y=2x-2$
\task $y=\dfrac{1}{2}x-1$
\task $y=-\dfrac{1}{2}x+1$
\end{tasks}
\dapan{A}
\end{cauhoi}
\begin{cauhoi}{Câu 47}{9H1-C1-B1-D02-TN-TH-047}
Viết nghiệm tổng quát của phương trình $-3x+5y=-7$ (biểu diễn $y$ theo $x$).
\begin{tasks}(2)
\task $y=\dfrac{3}{5}x+\dfrac{7}{5}$
\task $y=\dfrac{5}{3}x+\dfrac{7}{3}$
\task $y=\dfrac{3}{5}x-\dfrac{7}{5}$
\task $y=-\dfrac{3}{5}x-\dfrac{7}{5}$
\end{tasks}
\dapan{C}
\end{cauhoi}
\begin{cauhoi}{Câu 48}{9H1-C1-B1-D02-TN-TH-048}
Viết nghiệm tổng quát của phương trình $3x+2y=5$ (biểu diễn $y$ theo $x$).
\begin{tasks}(2)
\task $y=-\dfrac{3}{2}x+\dfrac{5}{2}$
\task $y=\dfrac{3}{2}x+\dfrac{5}{2}$
\task $y=-\dfrac{2}{3}x+\dfrac{5}{3}$
\task $y=-\dfrac{3}{2}x-\dfrac{5}{2}$
\end{tasks}
\dapan{A}
\end{cauhoi}
\begin{cauhoi}{Câu 49}{9H1-C1-B1-D02-TN-TH-049}
Viết nghiệm tổng quát của phương trình $5x+4y=3$ (biểu diễn $y$ theo $x$).
\begin{tasks}(2)
\task $y=\dfrac{5}{4}x+\dfrac{3}{4}$
\task $y=-\dfrac{4}{5}x+\dfrac{3}{5}$
\task $y=-\dfrac{5}{4}x+\dfrac{3}{4}$
\task $y=-\dfrac{5}{4}x-\dfrac{3}{4}$
\end{tasks}
\dapan{C}
\end{cauhoi}
\begin{cauhoi}{Câu 50}{9H1-C1-B1-D02-TN-TH-050}
Viết nghiệm tổng quát của phương trình $2x+4y=2$ (biểu diễn $y$ theo $x$).
\begin{tasks}(2)
\task $y=-\dfrac{1}{2}x-\dfrac{1}{2}$
\task $y=\dfrac{1}{2}x+\dfrac{1}{2}$
\task $y=-\dfrac{1}{2}x+\dfrac{1}{2}$
\task $y=-2x+1$
\end{tasks}
\dapan{C}
\end{cauhoi}
\clearpage\section*{Bảng đáp án trắc nghiệm}
\begin{center}\begin{tabular}{|c|c|c|c|c|c|c|c|c|c|}
\hline
1. C & 2. A & 3. B & 4. C & 5. B & 6. C & 7. D & 8. B & 9. C & 10. B\\\hline
11. D & 12. B & 13. B & 14. B & 15. B & 16. A & 17. C & 18. D & 19. D & 20. A\\\hline
21. C & 22. D & 23. D & 24. A & 25. B & 26. D & 27. C & 28. C & 29. C & 30. B\\\hline
31. B & 32. D & 33. D & 34. D & 35. B & 36. B & 37. C & 38. C & 39. B & 40. C\\\hline
41. B & 42. B & 43. A & 44. C & 45. D & 46. A & 47. C & 48. A & 49. C & 50. C\\\hline
\end{tabular}\end{center}
\end{document}
